\section{Phase 2 --- what kind of structure? Stationarity, order, and scope}
\label{sec:phase2}

\subsection{H2 fails: the kernel is pervasively depth-dependent}
\label{sec:phase2:stationarity}

H2 asks whether the AR kernel is the \emph{same} kernel at every depth. The
pre-registered test partitions depth into thirds (early $\ell$ 0--7, middle
8--15, late 16--23), fits one band-shared AR($k$) kernel per band (per-depth
intercepts $\mu_\ell$ absorb mean offsets, so only transition structure is
compared; $k=5$, the plan's early-band feasibility cap
$\lfloor L/3 \rfloor - 3$; working rank 256 after the plan's $r$-vs-$2r$
spot-check promoted 128), and compares the kernels four ways.

Every comparison rejects stationarity (\cref{fig:h2}):
\emph{cross-band prediction} --- within-band held-out $\Rsq$ is positive
(0.74 / 0.58 / 0.67) but \textbf{every} off-diagonal is \textbf{negative}
(e.g.\ early$\to$middle $-1.00$, late$\to$middle $-1.89$): a kernel from one
band predicts another band's updates \emph{worse than that band's own mean}.
This holds at every truncation rank from 32 to 256 (worst off-diagonal
$-0.32$ even at rank 32). \emph{Frobenius distance} between band kernels is
1.39--1.42 $\approx \sqrt{2}$, the orthogonal-of-equal-norm ceiling;
\emph{principal angles} between kernel output subspaces average
57--62$^\circ$; and a \emph{bootstrap} within bands puts the between-band
distance \textbf{80--104 standard deviations} above within-band sampling
spread. Crucially there is \emph{no clean phase boundary}: the three pairwise
distances are equal to within 3\%, so adjacent bands differ as much as distant
ones --- the depth dependence is \textbf{pervasive drift}, not segmentation.

Three caveats travel with this result: the cap makes it a short-horizon test
(long-horizon stationarity is rechecked, and also fails, on the fitted
generative model); AR-on-$\Delta$ stationarity is a proxy for H2's generative
(MA) form --- but a proxy whose failure is decisive in the rejecting
direction; and the $\pm\vl$ shared term biases the comparison \emph{toward}
stationarity, so rejection is the conservative outcome.

\begin{figure}[t]
  \centering
  \includesvg[width=0.62\linewidth]{fig04_h2_cross_band}
  \caption{H2 stationarity test on \pythia{}: held-out $\Rsq$ of band-shared
  AR(5) kernels evaluated within and across depth bands. Positive diagonal,
  every off-diagonal negative --- a kernel from one depth band is worse than
  the target band's mean. Rank-robust; bands bootstrap-separated by 80--104
  SD; no segment boundary.}
  \label{fig:h2}
\end{figure}

\textbf{Gate-2 verdict (pre-registered decision matrix): H2 fails
pervasively $\Rightarrow$ Phase 3 fits the matrix's state-space branch} --- a
state-space model with
\emph{free per-layer operators} (the smoothly-varying variant assumes exactly
the smoothness the data reject), and --- per the proposal's own framing ---
CLT-style unconstrained per-(source,\,target) parameterisation is
\emph{validated}: ``shared kernel along depth'' is the wrong abstraction for
these dynamics. The parsimony claim originally attached to H2 is off; what
survives is a reconstruction comparison at stated parameter budgets
(\cref{sec:evaluation:clt}).

\subsection{Effective order: \texorpdfstring{$\kstar \approx 1$}{k* approx 1},
everywhere we looked}
\label{sec:phase2:kstar}

Per-depth AR($k$) sweeps define the effective depth $\kstar$ as the smallest
$k$ within 95\% of the asymptotic $\Rsq$. On the factual store
$\boldsymbol{\kstar \approx 1}$ at every depth (median 1; the single
exception, $\kstar{=}2$ at $\ell{=}17$, coincides with the answer-promotion
write discussed in \cref{sec:evaluation:probe}): one preceding update captures
essentially all held-out predictable variance (per-depth best $\Rsq$
0.83--0.94), lags $\ge 2$ add a median $+0.018$, and the scrambled null
collapses the whole curve to 0 (\cref{fig:kstar}). The contrast with the
Granger reach of 20 (\cref{fig:h1-granger}) is a headline in itself:
\textbf{long-range depth structure is statistically real but
reconstruction-negligible} --- and \cref{sec:phase1:sublayer} locates the two
halves of that sentence in different sublayers (deep reach: attention;
reconstruction mass: MLP). $\kstar$ inherits the lag-1 tautology caveat, but
it is also exactly the model-order information Phase 3 needs, and the fitted
model later supplies a mechanism for it (a contractive state operator,
\cref{sec:evaluation:operator}). Stratifying by relation class leaves
$\kstar=1$ in every class and every depth band (class-invariant order); on
generic text the median is $\kstar=2$ (band pattern 1/3/1) --- stated
explicitly since it differs from the factual 1, with per-depth generic fits
being visibly noisier.

\begin{figure}[t]
  \centering
  \includesvg[width=0.78\linewidth]{fig05_kstar_vs_depth}
  \caption{Per-depth effective order $\kstar$ on \pythia{} (bars), with the
  per-depth lag-1 $\Rsq$ (purple line, right axis) and the Granger reach (20)
  for contrast. The
  one $\kstar{=}2$ depth, $\ell{=}17$, is the answer-promotion write: the state
  probe jumps there (\cref{fig:probe}) and the source-prefix sweep dips there
  (\cref{fig:prefix}).}
  \label{fig:kstar}
\end{figure}

\subsection{Scope: the kernel is query-position-specific and circuit-specific}
\label{sec:phase2:scope}

\paragraph{Token position.}
Re-collecting trajectories at four positions per prompt (first subject token,
last subject token, a mid-prompt non-subject token, final prompt token) and
refitting AR($\kstar{=}1$) per position: within-position held-out $\Rsq$ is
positive at every position (0.28--0.55), but \textbf{every cross-position
$\Rsq$ is negative} (mean $-0.36$; kernels near-orthogonal, Frobenius
$\approx\sqrt{2}$, angles 57--60$^\circ$) --- the same no-transfer signature
as the depth bands, now across positions. The scrambled null collapses the
diagonals, and the tautology (shared by diagonal and off-diagonal) cancels in
the retention ratio. The Phase-1/2 characterisation is therefore scoped to
\emph{context-aggregated factual recall at the query position}; Phase 3 fits
there, and every claim in this paper inherits that scope.

\paragraph{Relation class.}
The mirror-image test across the three relation classes (run with the fitted
SSM of \cref{sec:decomposition}; reported here because it completes the
scope picture): within-class held-out $\Rsq$ mean $+0.835$, \textbf{every
cross-class transfer $\approx 0$ or negative} (mean $-0.191$;
\cref{fig:crossclass}). The dynamics are finely tuned to each recall circuit
--- \emph{circuit-specific kernels} --- while the \emph{order} $\kstar=1$ is
class-invariant: same model order everywhere, different kernels per circuit.

\begin{figure}[t]
  \centering
  \includesvg[width=0.6\linewidth]{fig10_crossclass}
  \caption{Cross-class transfer of the fitted decomposition (train class
  $\to$ eval class, held-out $\Rsq$). Strong diagonal, no transfer off it ---
  the same signature as cross-band and cross-position: structured, but
  specific.}
  \label{fig:crossclass}
\end{figure}

\subsection{Cross-model replication: six models, three families}
\label{sec:phase2:crossmodel}

The full H1/H2 pipeline (and, where public residual SAEs exist, the Phase-3
fit) was replicated on five further models; Pythia-70m additionally serves as
the optimiser-debug lane and a depth point in
\cref{sec:evaluation:depth}. \Cref{tab:crossmodel} shows the pattern is
universal across the models tested: AR-on-$\Delta$ predictability far above
its null with full Granger significance, every cross-band $\Rsq$ negative,
and $\kstar = 1$ at every band of every model.

\begin{table}[t]
  \centering
  \small
  \caption{Cross-model replication of H1/H2. AR $\Rsq$ is pooled held-out
  real$-$null; Granger counts are Bonferroni-significant cells (real vs
  scrambled null) out of all feasible (depth, lag) cells --- $\sum_j (L - j)$
  over the lag grid: $\{1,2,3,5,10,15,20\}$ for $L \ge 24$ (112 cells at
  $L{=}24$, 196 at $L{=}36$); the 12-layer models use the capped grid
  $\{1,2,3,5,8,10\}$ since lags 15 and 20 are infeasible at $L{=}12$, giving
  43 cells (\cref{app:stats}). H2 column reports the sign pattern of
  cross-band transfer. Datasets are per-model confidence-filtered CounterFact
  (sizes 767--2{,}540).}
  \label{tab:crossmodel}
  \begin{tabular}{llcccc}
    \toprule
    Model & $L$ & AR $\Rsq$ (real$-$null) & Granger real / null &
    cross-band $\Rsq$ & $\kstar$ \\
    \midrule
    GPT-2-small           & 12 & 0.625 & 43/43 \;/\; 0 & all $<0$ & 1 \\
    Pythia-160m-deduped   & 12 & 0.894 & 43/43 \;/\; 0 & all $<0$ & 1 \\
    GPT-2-medium          & 24 & 0.513 & 112/112 \;/\; 0 & all $<0$ & 1 \\
    Qwen2-0.5B            & 24 & 0.462 & 112/112 \;/\; 0 & all $<0$ & 1 \\
    \pythia{} (primary)   & 24 & 0.650 & 112/112 \;/\; 0 & all $<0$ & 1 \\
    GPT-2-large           & 36 & 0.311 & 196/196 \;/\; 0 & all $<0$ & 1 \\
    \bottomrule
  \end{tabular}
\end{table}

One anomaly is recorded rather than hidden: on GPT-2-small the cheap
\emph{low-rank} ridge diagnostic scores only 0.146 layer-averaged despite a
full-rank 0.91 and AR 0.625 --- GPT-2's $v\!\to\!\Delta$ map is far less
low-rank than Pythia's. It does not affect any H1/H2 verdict.

\subsection{Generality beyond factual recall}
\label{sec:phase2:generic}

Re-running the identical pipelines on 1{,}555 generic Pile windows ---
length-matched to the factual prompt multiset (content lengths 4--24 tokens,
median 9; one window per distinct document; this narrow-window construction is
a stated scope condition) --- replicates the \emph{entire} characterisation:
H1's binary claim holds (null $\approx 0$, Granger 112/112 vs 0/112,
lag$\ge$2 89/89), H2 fails even more emphatically (every cross-band $\Rsq$
negative, worst $-144$; bands separated by 39--188 bootstrap SD), and the
effective order is low ($\kstar$ median 2). The raw generic AR level (0.935)
is \emph{higher} than factual (0.650), and \cref{sec:phase1:pressure} is the
reason that must not be read as ``more computation'': T5 shows the generic
level is $\approx 0.94$ tautological $+\;{\approx}0$ genuine, and after
persistence-stripping generic structure (0.29) is \emph{weaker} than factual
(0.44). The phenomenon is general to transformer depth; the \emph{kernels}
remain position- and circuit-specific per \cref{sec:phase2:scope}.
